%%%%%%%%%%%%%%%%%%%%%%%
\section{Introduction}
%%%%%%%%%%%%%%%%%%%%%%%
Chisel (\textbf{C}onstructing \textbf{H}ardware \textbf{I}n a \textbf{S}cala \textbf{E}mbedded \textbf{L}anguage) is a hardware design language embedded within Scala. It offers hardware designers the ability to write complex, parameterizable, and reusable hardware components using high-level programming constructs \cite{bachrach2012chisel}. 

This Appendix is a very brief introduction to Chisel primarily for designers to quickly reference features. For processor design, \citep{chisel:book} is a highly recommended free book.

%%%%%%%%%%%%%%%%%%%%%%%%%%%%
\section{Overview of Chisel}
%%%%%%%%%%%%%%%%%%%%%%%%%%%%
Chisel is a domain-specific language (DSL) for hardware design, leveraging Scala's object-oriented and functional programming features \cite{chisel-doc}. It compiles to synthesizable Verilog, integrating with traditional hardware development flows.

\subsection{Key Features}

\begin{itemize}
    \item Hardware Construction Language: High-level abstractions for hardware components.
    \item Parameterization: Easy creation of generic and reusable modules.
    \item Strong Typing: Reduces errors through rigorous type checking.
    \item Functional Programming: Utilizes Scala's functional paradigms for concise code.
\end{itemize}

\section{Basic Syntax}

\subsection{Scala Basics}

Understanding basic Scala syntax is essential for Chisel programming \cite{odersky2008programming}.

\begin{itemize}
    \item Immutable Variables: Declared with \texttt{val}, cannot be reassigned.
    \item Mutable Variables: Declared with \texttt{var}, can be reassigned.
    \item Functions: Defined using the \texttt{def} keyword.
    \item Classes and Objects: Used to define hardware modules and singletons.
\end{itemize}

\subsection{Import Statements}

Include necessary Chisel libraries at the beginning of your file:


%\begin{lstlisting}[language=Scala]
{
\begin{chiselcode}
import chisel3._
import chisel3.util._    
\end{chiselcode}
}
%\end{lstlisting}

\section{Data Types}

Chisel provides specialized data types for hardware representation \cite{chisel-doc}.

\subsection{Base Types}

\begin{itemize}
    \item Bool: Single-bit boolean (\texttt{true}/\texttt{false}).
        \newline Declaration: \texttt{val myBool = Bool()}
    \item UInt: Unsigned integer with specified width.
        \newline Declaration: \texttt{val myUInt = UInt(8.W)}
    \item SInt: Signed integer with specified width.
        \newline Declaration: \texttt{val mySInt = SInt(16.W)}
    \item FixedPoint: Fixed-point number with total and fractional bits.
        \newline Declaration: \texttt{val myFixed = FixedPoint(16.W, 8.BP)}
    \item Clock: Represents a clock signal.
    \item Reset: Represents a reset signal.
\end{itemize}

\subsection{Aggregate Types}

\begin{itemize}
    \item Bundle: Custom data structures combining multiple fields.
        \newline Declaration: \texttt{class MyBundle extends Bundle \{ val a = UInt(8.W) \}}
    \item Vec: Indexable collection of elements of the same type.
        \newline Declaration: \texttt{val myVec = Vec(4, UInt(8.W))}
\end{itemize}

\subsection{Literals}

\begin{itemize}
    \item Unsigned Integer: \texttt{10.U}, \texttt{3.U(8.W)}
    \item Signed Integer: \texttt{-5.S}, \texttt{-2.S(8.W)}
    \item Boolean: \texttt{true.B}, \texttt{false.B}
\end{itemize}

\subsection{Operators}

\begin{itemize}
    \item Arithmetic: \texttt{+}, \texttt{-}, \texttt{*}, \texttt{/}, \texttt{\%}
    \item Logical: \texttt{\&}, \texttt{|}, \texttt{\^}, \texttt{!}
    \item Comparison: \texttt{===}, \texttt{=/=}, \texttt{<}, \texttt{<=}, \texttt{>}, \texttt{>=}
    \item Bitwise Reduction: \texttt{\&}, \texttt{|}, \texttt{\^} (applied to a single operand)
    \item Shift: \texttt{<<}, \texttt{>>}
\end{itemize}

\section{Modules and Hierarchy}

\subsection{Defining Modules}

Modules are defined by extending the \texttt{Module} class \cite{chisel-doc}.

\begin{lstlisting}[language=Scala]
class MyModule extends Module {
  // Module implementation
}
\end{lstlisting}

\subsection{Input/Output Ports}

Ports are declared using the \texttt{IO} method with a \texttt{Bundle}.

\begin{lstlisting}[language=Scala]
class MyModule extends Module {
  val io = IO(new Bundle {
    val in  = Input(UInt(8.W))
    val out = Output(UInt(8.W))
  })
  // Module logic
}
\end{lstlisting}

\subsection{Instantiating Modules}

Modules can be instantiated within other modules.

\begin{lstlisting}[language=Scala]
val subModule = Module(new SubModule)
subModule.io.in := io.in
io.out := subModule.io.out
\end{lstlisting}

\subsection{Parameterized Modules}

Modules can be parameterized for reusability.

\begin{lstlisting}[language=Scala]
class MyModule(val width: Int) extends Module {
  val io = IO(new Bundle {
    val in  = Input(UInt(width.W))
    val out = Output(UInt(width.W))
  })
  // Module logic
}
\end{lstlisting}

\section{Wires and Registers}

\subsection{Wires}

Wires represent combinational logic signals \cite{chisel-doc}.

\begin{lstlisting}[language=Scala]
val tempWire = Wire(UInt(8.W))
tempWire := io.in + 1.U
\end{lstlisting}

\subsection{Registers}

Registers store state across clock cycles.

\begin{itemize}
    \item \texttt{Reg}: Creates a register without an initial value.
        \newline Declaration: \texttt{val myReg = Reg(UInt(8.W))}
    \item \texttt{RegInit}: Creates a register with an initial value.
        \newline Declaration: \texttt{val myReg = RegInit(0.U(8.W))}
\end{itemize}

\subsection{Register Updates}

Registers are updated using non-blocking assignment.

\begin{lstlisting}[language=Scala]
myReg := myReg + 1.U
\end{lstlisting}

\section{Control Structures}

\subsection{\texttt{when} Statements}

Conditional hardware execution.

\begin{lstlisting}[language=Scala]
when (condition) {
  // Actions when condition is true
} .elsewhen (anotherCondition) {
  // Actions when anotherCondition is true
} .otherwise {
  // Actions when all conditions are false
}
\end{lstlisting}

\subsection{\texttt{switch} Statements}

Multi-way conditional execution.

\begin{lstlisting}[language=Scala]
switch(selector) {
  is(value1) {
    // Actions for value1
  }
  is(value2) {
    // Actions for value2
  }
}
\end{lstlisting}

\subsection{\texttt{Mux} Function}

Multiplexer for conditional selection.

\begin{lstlisting}[language=Scala]
val result = Mux(condition, trueValue, falseValue)
\end{lstlisting}

\subsection{Loops}

Used for hardware generation, not for creating sequential logic \cite{chisel-tutorial}.

\begin{lstlisting}[language=Scala]
for (i <- 0 until n) {
  // Generate n instances of hardware
}
\end{lstlisting}

\section{Functions and Procedures}

\subsection{Defining Functions}

Functions encapsulate reusable logic.

\begin{lstlisting}[language=Scala]
def add(a: UInt, b: UInt): UInt = {
  a + b
}
\end{lstlisting}

\subsection{Invoking Functions}

\begin{lstlisting}[language=Scala]
val sum = add(io.in1, io.in2)
\end{lstlisting}

\section{Parameters and Generics}

\subsection{Parameterized Types}

Types can be parameterized for flexibility.

\begin{lstlisting}[language=Scala]
def myFunction[T <: Data](in: T): T = {
  // Function body
}
\end{lstlisting}

\subsection{Generic Modules}

Modules can be generic over data types.

\begin{lstlisting}[language=Scala]
class MyGenericModule[T <: Data](genType: T) extends Module {
  val io = IO(new Bundle {
    val in  = Input(genType)
    val out = Output(genType)
  })
  // Module logic
}
\end{lstlisting}

\section{Hardware Construction in Chisel}

\subsection{Vectors (\texttt{Vec})}

Indexable collections of elements.

\begin{lstlisting}[language=Scala]
val myVec = VecInit(Seq.fill(4)(0.U(8.W)))
myVec(0) := io.in
\end{lstlisting}

\subsection{Memories (\texttt{Mem})}

Memory elements for storage.

\begin{lstlisting}[language=Scala]
val myMem = Mem(1024, UInt(8.W))
val dataOut = myMem.read(addr)
myMem.write(addr, dataIn)
\end{lstlisting}

\subsection{Bundles}

Custom data structures grouping multiple fields.

\begin{lstlisting}[language=Scala]
class MyBundle extends Bundle {
  val field1 = UInt(8.W)
  val field2 = Bool()
}
\end{lstlisting}

\section{Testing and Verification}

\subsection{ChiselTest Framework}

ChiselTest provides tools for testing Chisel modules \cite{chiseltest}.

\begin{lstlisting}[language=Scala]
import chiseltest._
import org.scalatest.flatspec.AnyFlatSpec

class MyModuleTest extends AnyFlatSpec with ChiselScalatestTester {
  "MyModule" should "perform correctly" in {
    test(new MyModule) { c =>
      c.io.in.poke(0.U)
      c.clock.step(1)
      c.io.out.expect(0.U)
    }
  }
}
\end{lstlisting}

\subsection{Assertions}

In-module checks using \texttt{assert}.

\begin{lstlisting}[language=Scala]
assert(io.out === expectedValue)
\end{lstlisting}

\section{Peculiarities of Chisel}

\subsection{Chisel is Embedded in Scala}

Chisel code is Scala code. This means:

\begin{itemize}
    \item Full access to Scala's features.
    \item Potential confusion between hardware constructs and software constructs.
\end{itemize}

\subsection{Delayed Execution Model}

Chisel builds a hardware graph during elaboration, which may lead to unexpected behavior if side effects are not managed properly \cite{chisel-doc}.

\subsection{Assignment Operators}

\begin{itemize}
    \item \texttt{:=}: Used for updating hardware signals.
    \item \texttt{=}: Used for variable assignment in Scala.
\end{itemize}

\subsection{Type System}

Chisel enforces strict type checking, which may differ from SystemVerilog or VHDL.

\subsection{Combining Hardware and Software Constructs}

Be cautious when mixing Scala control structures (e.g., \texttt{if}, \texttt{for}) with Chisel constructs (\texttt{when}, \texttt{switch}) \cite{chisel-tutorial}.

\subsection{Hardware vs. Software Mindset}

Designers must think in terms of hardware concurrency, even when writing code that resembles sequential software.

\section{Conclusion}

Chisel provides a modern approach to hardware design, combining the rigor of hardware description with the flexibility of a high-level programming language. By understanding Chisel's syntax and peculiarities, hardware designers can leverage its powerful features to create complex and reusable hardware components.